\section{Fourier 分析：从三角矩到热核}\label{sec:07-02}

\subsection{圆周 Fourier 级数、\Fejer{} 逼近与 Plancherel}\label{subsec:7-2-1}
\index{圆周 Fourier 级数}\index{Fejer@\Fejer{} 核}\index{Plancherel 定理}

把圆周写成 $\T=\R/2\pi\Z$，并用
\[
  \dd m(x)=\frac{\dd x}{2\pi}
\]
表示归一化 Haar 测度。指数函数 $e_n(x)=e^{inx}$ 满足
$\int_\T e_n\overline{e_k}\,\dd m=\1_{\{n=k\}}$。因此，若只观测函数对所有
$e_n$ 的积分，所得数列正是函数关于这族正交方向的坐标。困难不在于坐标是否存在，而在于这些坐标
能否把原函数恢复出来。

最先想到的恢复算子是对称部分和
\[
  S_Nf(x)=\sum_{|n|\leq N}\widehat f(n)e^{inx}.
\]
下一个例子说明，不能把有限维正交投影的直觉直接升级为逐点恢复。

\begin{example}[Gibbs 现象与 Dirichlet 部分和的失效]\label{ex:7-2-1-1}
令 $f$ 是圆周上的方波，取代表
\[
 f(x)=
 \begin{cases}
   -1,&-\pi<x<0,\\
   0,&x=0\text{ 或 }x=\pi\equiv-\pi,\\
   1,&0<x<\pi.
 \end{cases}
\]
直接积分得到：当 $n\neq0$ 时
\[
 \widehat f(n)
 =\frac1{2\pi}\left(\int_0^\pi-\int_{-\pi}^0\right)e^{-inx}\dd x
 =\frac{1-(-1)^n}{\pi i n}.
\]
于是，以奇数频率 $2N+1$ 截断时，
\begin{equation}\label{eq:7-2-gibbs-partial-sum}
 S_{2N+1}f(x)=\frac4\pi\sum_{k=0}^N\frac{\sin((2k+1)x)}{2k+1}.
\end{equation}
对右端求导，并使用有限等比和，得到
\[
 \frac{\dd}{\dd x}S_{2N+1}f(x)
 =\frac4\pi\sum_{k=0}^N\cos((2k+1)x)
 =\frac2\pi\frac{\sin(2(N+1)x)}{\sin x}.
\]
取 $x_N=\pi/(2(N+1))$，由 $S_{2N+1}f(0)=0$ 及换元
$u=2(N+1)x$，有
\[
 S_{2N+1}f(x_N)
 =\frac2\pi\int_0^\pi
   \frac{\sin u}{2(N+1)\sin(u/(2(N+1)))}\dd u
 \longrightarrow \frac2\pi\int_0^\pi\frac{\sin u}{u}\dd u.
\]
这里的支配可取常数倍的 $\sin u/u$：在 $0\leq v\leq\pi/2$ 上
$\sin v\geq 2v/\pi$。极限严格大于 $1$。例如把
$\sin u/u$ 展开并使用交错级数余项，
\[
 \int_0^\pi\frac{\sin u}{u}\dd u
 >\pi-\frac{\pi^3}{18}+\frac{\pi^5}{600}
       -\frac{\pi^7}{35280}>\frac\pi2.
\]
数值上该极限约为 $1.17898$。因此观察点虽趋向跳点，过冲却不趋于零；这就是 Gibbs 现象。

这种失效不只来自跳跃。Dirichlet 核
\[
 D_N(x)=\sum_{|n|\leq N}e^{inx}
       =\frac{\sin((N+\tfrac12)x)}{\sin(x/2)}
\]
的 $L^1(m)$ 范数按 $\log N$ 增长。下面给出带显式区间的下界。对
$1\leq k\leq\lfloor N/2\rfloor$，置
\[
 I_{k,N}=\left[
   \frac{(k+1/3)\pi}{N+1/2},
   \frac{(k+2/3)\pi}{N+1/2}
 \right].
\]
这些区间两两不交。若 $x\in I_{k,N}$，则
\[
 \left|\sin\bigl((N+\tfrac12)x\bigr)\right|\geq\frac{\sqrt3}{2},
 \qquad 0<\sin(x/2)\leq\frac{x}{2},
\]
因而 $|D_N(x)|\geq\sqrt3/x$。使用归一化测度 $\dd m=\dd x/(2\pi)$，逐段得到
\begin{align*}
 \int_{I_{k,N}}|D_N(x)|\dd m(x)
 &\geq\frac{\sqrt3}{2\pi}
       \frac{|I_{k,N}|}{\sup I_{k,N}}\\
 &=\frac{\sqrt3}{6\pi(k+2/3)}
 \geq\frac{\sqrt3}{10\pi}\frac1k.
\end{align*}
因此
\[
 \|D_N\|_{L^1(m)}
 \geq\frac{\sqrt3}{10\pi}
       \sum_{k=1}^{\lfloor N/2\rfloor}\frac1k
 \geq c'\log N
\]
（$N$ 充分大），其中 $c'>0$ 与 $N$ 无关。

令 $L_N:C(\T)\to\C$ 为 $L_Ng=S_Ng(0)=\int g(-y)D_N(y)\dd m(y)$。
三角不等式先给 $\|L_N\|\leq\|D_N\|_1$。反向固定 $\eta>0$。$D_N$ 只有有限多个零点，可在这些
零点周围取有限个小弧，其并记为 $O$，并使
$\int_O|D_N|\dd m<\eta$。在补集上令
$g(-y)=\operatorname{sgn}D_N(y)$，再在每个小弧上作模不超过 $1$ 的线性连接，得到
$g\in C(\T)$、$\|g\|_\infty\leq1$。于是
\[
 |L_Ng|
 \geq\int_{O^c}|D_N|\dd m-\int_O|D_N|\dd m
 \geq\|D_N\|_1-2\eta.
\]
令 $\eta\downarrow0$，得到 $\|L_N\|=\|D_N\|_1$。若每个 $g\in C(\T)$ 的数列
$(L_Ng)$ 都有界，则闭集
\[
 E_r=\{g:\sup_N|L_Ng|\leq r\}\qquad(r\in\N)
\]
覆盖完备空间 $C(\T)$。Baire 定理给出某个球
$B(g_0,\rho)\subset E_r$。对 $\|h\|_\infty\leq1$，把
$g_0+(\rho/2)h$ 与 $g_0$ 相减，便得
$\sup_N|L_Nh|\leq4r/\rho$，这与 $\|L_N\|\to\infty$ 矛盾。
所以存在连续函数 $g$，其 Fourier 部分和甚至在固定点 $0$ 发散。
\end{example}

Dirichlet 核的问题是绝对质量越来越大而且正负振荡。\Cesaro{} 平均把前 $N+1$ 个部分和再平均；
这看似只是数列求和法的改变，却恰好把振荡核换成了正核。

\begin{example}[\Fejer{} 正核]\label{ex:7-2-1-2}
有限和的次数可以交换，因此
\begin{align*}
 F_N(x)
 &=\frac1{N+1}\sum_{k=0}^ND_k(x)\\
 &=\sum_{|n|\leq N}\left(1-\frac{|n|}{N+1}\right)e^{inx}\\
 &=\frac1{N+1}\left|\sum_{k=0}^Ne^{ikx}\right|^2
 =\frac1{N+1}\left(\frac{\sin((N+1)x/2)}{\sin(x/2)}\right)^2.
\end{align*}
在 $x\in2\pi\Z$ 处取连续延拓值 $N+1$。平方表示说明 $F_N\geq0$；常数 Fourier
系数为 $1$，故 $\int_\T F_N\dd m=1$。固定 $0<\delta<\pi$ 后，若
$\delta\leq|x|\leq\pi$，则
\[
 F_N(x)\leq\frac1{(N+1)\sin^2(\delta/2)},
 \qquad
 \int_{\delta\leq|x|\leq\pi}F_N(x)\dd m(x)\longrightarrow0.
\]
所以全部质量集中到圆周原点。并且若 $f\geq0$ a.e.，则
$f*F_N\geq0$；这项保正性正是 Dirichlet 部分和所没有的稳定性。
\end{example}

Fourier 系数也可以对测度定义。此时系数的高频行为能排除某些对象，却不能单独确定测度的全部
几何性质。

\begin{example}[原子测度]\label{ex:7-2-1-3}
对 $x_0\in\T$，
\[
 \widehat{\delta_{x_0}}(n)
 =\int_\T e^{-inx}\dd\delta_{x_0}(x)=e^{-inx_0},
 \qquad |\widehat{\delta_{x_0}}(n)|=1.
\]
因此原子测度的 Fourier 系数不趋零。推论~\ref{corollary:7-2-1-2} 与圆周上的
Riemann--Lebesgue 结论还将说明：这种测度不可能有 $L^1(m)$ 密度。

反过来，“系数不趋零”并不推出存在原子。令
\[
 X=2\pi\sum_{j=1}^\infty\frac{2\varepsilon_j}{3^j},
 \qquad \Prob(\varepsilon_j=0)=\Prob(\varepsilon_j=1)=\frac12,
\]
并令 $\mu$ 为 $X$ 在圆周上的分布。这是无原子的 Cantor 型奇异测度，而独立性给出
\[
 |\widehat\mu(n)|
 =\prod_{j=1}^\infty\left|\cos\frac{2\pi n}{3^j}\right|.
\]
取 $n=3^k$，前 $k$ 个因子绝对值为 $1$，故
\[
 |\widehat\mu(3^k)|
 =\prod_{r=1}^\infty\left|\cos\frac{2\pi}{3^r}\right|>0.
\]
先核对“无原子”。一个 Cantor 点至多对应两个只含数字 $0,2$ 的三进制展开；任一指定数字串的前
$N$ 位同时出现的概率为 $2^{-N}$，故该完整数字串的概率至多
$\lim_{N\to\infty}2^{-N}=0$。因此对每个 $x$，
$\mu(\{x\})\leq2\lim_N2^{-N}=0$。

最后的无穷乘积为正也可逐项核验。令 $\theta_r=2\pi/3^r$。当 $r\geq2$ 时
$0<\cos\theta_r<1$，并且
\[
 u_r:=1-\cos\theta_r\leq\frac{\theta_r^2}{2},
 \qquad \sum_{r=2}^\infty u_r<\infty.
\]
又有 $u_r<1/2$，所以 $-\log(1-u_r)\leq2u_r$；从而
$\sum_{r\geq2}-\log(\cos\theta_r)<\infty$，即
$\prod_{r\geq2}\cos\theta_r>0$。乘上第一个因子的绝对值
$|\cos(2\pi/3)|=1/2$，便得所写无穷乘积严格为正。这给出非原子而 Fourier 系数不衰减的具体例子。
\end{example}

\begin{definition}[Fourier 系数]\label{def:7-2-1-1}
若 $f\in L^1(\T,m)$，定义
\[
 \widehat f(n)=\int_\T f(x)e^{-inx}\dd m(x),\qquad n\in\Z.
\]
若 $\mu$ 是 $\T$ 上有限复 Borel 测度，则定义
$\widehat\mu(n)=\int_\T e^{-inx}\dd\mu(x)$。对可积函数，圆周卷积约定为
\[
 (f*g)(x)=\int_\T f(x-y)g(y)\dd m(y).
\]
在 $L^2(\T)$ 中取内积 $\langle f,g\rangle=\int f\overline g\dd m$，于是
$\widehat f(n)=\langle f,e_n\rangle$。
\end{definition}

\begin{definition}[Dirichlet 与 \Fejer{} 核]\label{def:7-2-1-2}
对 $N\geq0$，定义
\[
 D_N=\sum_{|n|\leq N}e_n,
 \qquad
 F_N=\frac1{N+1}\sum_{k=0}^ND_k.
\]
相应的 Fourier 部分和与 \Fejer{} 平均为
\[
 S_Nf=f*D_N,
 \qquad
 \sigma_Nf=\frac1{N+1}\sum_{k=0}^NS_kf=f*F_N.
\]
\end{definition}

例~\ref{ex:7-2-1-2} 已证明 $F_N\geq0$、$\int F_N\dd m=1$，且其远离零点的质量趋于零；
即 $(F_N)$ 是圆周上的正近似恒等族。

\begin{definition}[三角矩问题]\label{def:7-2-1-3}
给定复数列 $(c_n)_{n\in\Z}$，三角矩问题是寻找函数 $f$ 或有限测度 $\mu$，使
\[
 c_n=\int_\T e^{-inx}f(x)\dd m(x)
 \quad\text{或}\quad
 c_n=\int_\T e^{-inx}\dd\mu(x)
 \qquad(n\in\Z).
\]
\end{definition}

\begin{proposition}[正定性的必要性]\label{proposition:7-2-1-positive-definite-necessity}
若正测度 $\mu$ 解三角矩问题，则对任意 $r\geq1$、整数 $n_1,\ldots,n_r$ 与复数
$z_1,\ldots,z_r$，都有
\[
 \sum_{j,k=1}^r c_{n_j-n_k}z_j\overline{z_k}\geq0.
\]
\end{proposition}

\begin{proof}
把矩表示代入有限和，得到
\[
 \sum_{j,k=1}^r c_{n_j-n_k}z_j\overline{z_k}
 =\int_\T\left|\sum_{j=1}^rz_je^{-in_jx}\right|^2\dd\mu(x)\geq0,
\]
因为被积函数与测度均为非负。
\end{proof}

因此正定性是正测度三角矩的必要条件。第~\ref{subsec:7-3-1} 小节将证明相应的充分性；本小节先
解决 Fourier 系数所决定测度的唯一性。

\begin{theorem}[\Fejer{} 定理]\label{theorem:7-2-1-1}
若 $f\in C(\T)$，则 $\sigma_Nf\to f$ 一致收敛。若
$f\in L^p(\T)$ 且 $1\leq p<\infty$，则
$\sigma_Nf\to f$ 于 $L^p(\T)$。
\end{theorem}

\begin{proof}
先设 $f\in C(\T)$。由 $\int F_N\dd m=1$，
\begin{equation}\label{eq:7-2-fejer-difference}
 \sigma_Nf(x)-f(x)
 =\int_\T\bigl(f(x-y)-f(x)\bigr)F_N(y)\dd m(y).
\end{equation}
给定 $\varepsilon>0$，圆周紧性使 $f$ 一致连续；可取 $0<\delta<\pi$，使
$|f(x-y)-f(x)|<\varepsilon/2$ 对所有 $x$ 及 $|y|<\delta$ 成立。把积分分成近端与远端，并使用
$F_N\geq0$，得到
\begin{align*}
 \|\sigma_Nf-f\|_\infty
 &\leq \frac\varepsilon2\int_{|y|<\delta}F_N(y)\dd m(y)
 +2\|f\|_\infty\int_{|y|\geq\delta}F_N(y)\dd m(y)\\
 &\leq\frac\varepsilon2
 +\frac{2\|f\|_\infty}{(N+1)\sin^2(\delta/2)}.
\end{align*}
当 $N$ 足够大时第二项小于 $\varepsilon/2$，证明一致收敛。

再设 $f\in L^p$，$1\leq p<\infty$。积分型 Minkowski 不等式与平移等距性给出
\begin{equation}\label{eq:7-2-fejer-lp}
 \|\sigma_Nf-f\|_p
 \leq\int_\T\|f(\,\cdot-y)-f\|_pF_N(y)\dd m(y).
\end{equation}
这里平移在 $L^p(\T)$ 中强连续：先用连续函数 $g$ 逼近 $f$，再由
\[
 \|\tau_yf-f\|_p
 \leq2\|f-g\|_p+\|\tau_yg-g\|_p
 \leq2\|f-g\|_p+\|\tau_yg-g\|_\infty
\]
及 $g$ 的一致连续性令 $y\to0$。因此给定 $\varepsilon>0$，可取 $\delta>0$ 使
$\|f(\cdot-y)-f\|_p<\varepsilon/2$ 对 $|y|<\delta$ 成立。远端上该范数不超过
$2\|f\|_p$。于是与上面相同的分裂给出
\[
 \|\sigma_Nf-f\|_p
 \leq\frac\varepsilon2
 +2\|f\|_p\int_{|y|\geq\delta}F_N(y)\dd m(y)<\varepsilon
\]
，只要 $N$ 足够大。
\end{proof}

\begin{corollary}[Stone--Weierstrass 的 Fourier 版]\label{corollary:7-2-1-2}
三角多项式在 $C(\T)$ 中关于一致范数稠密，也在 $L^p(\T)$ 中关于 $L^p$ 范数稠密，
$1\leq p<\infty$。更具体地，$\sigma_Nf$ 就是一列只使用 $|n|\leq N$ 频率的逼近。
\end{corollary}

\begin{proof}
由例~\ref{ex:7-2-1-2} 的有限 Fourier 展开，
\[
 \sigma_Nf(x)
 =\sum_{|n|\leq N}\left(1-\frac{|n|}{N+1}\right)\widehat f(n)e^{inx},
\]
所以 $\sigma_Nf$ 是三角多项式。定理~\ref{theorem:7-2-1-1} 分别给出一致范数与
$L^p$ 范数收敛，因而给出两种稠密性。
\end{proof}

这里还得到圆周上的 Riemann--Lebesgue 结论。若 $f\in L^1(\T)$，给定
$\varepsilon>0$，由本推论取三角多项式 $p$ 使 $\|f-p\|_1<\varepsilon$。当 $|n|$ 超过
$p$ 的最高频率时，$\widehat p(n)=0$，所以
\[
 |\widehat f(n)|=|\widehat{f-p}(n)|\leq\|f-p\|_1<\varepsilon.
\]
因此 $\widehat f(n)\to0$ 当 $|n|\to\infty$。例~\ref{ex:7-2-1-3} 中的原子测度不可能有
$L^1$ 密度，正是这项结论的应用。

现在可以把 Bessel 的单向估计升级为 Parseval 等式。关键不是再作一次形式求和，而是利用刚刚得到的
稠密性证明指数族没有遗漏任何 $L^2$ 方向。

\begin{theorem}[$\T$ 上的 Plancherel/Parseval]\label{theorem:7-2-1-3}
Fourier 映射
\[
 \mathcal F_\T:L^2(\T)\longrightarrow\ell^2(\Z),
 \qquad f\longmapsto(\widehat f(n))_{n\in\Z},
\]
是满射线性等距同构。特别地，
\begin{equation}\label{eq:7-2-circle-parseval}
 \|f\|_2^2=\sum_{n\in\Z}|\widehat f(n)|^2,
 \qquad
 \langle f,g\rangle
 =\sum_{n\in\Z}\widehat f(n)\overline{\widehat g(n)}.
\end{equation}
并且对称部分和 $S_Nf=\sum_{|n|\leq N}\widehat f(n)e^{inx}$ 在 $L^2$ 中收敛到 $f$。
\end{theorem}

\begin{proof}
对有限集 $E\subset\Z$，令 $p_E=\sum_{n\in E}\widehat f(n)e_n$。正交性给出
\begin{align*}
 \|f-p_E\|_2^2
 &=\|f\|_2^2-2\operatorname{Re}\langle f,p_E\rangle+\|p_E\|_2^2\\
 &=\|f\|_2^2-\sum_{n\in E}|\widehat f(n)|^2.
\end{align*}
左边非负；令 $E$ 递增耗尽 $\Z$，得到 Bessel 不等式
$\sum_n|\widehat f(n)|^2\leq\|f\|_2^2$。

取 $E_N=\{-N,\ldots,N\}$。若 $q$ 是频率都落在 $E_N$ 中的三角多项式，则
$f-S_Nf$ 与 $q-S_Nf$ 正交，因而
\[
 \|f-q\|_2^2
 =\|f-S_Nf\|_2^2+\|S_Nf-q\|_2^2
 \geq\|f-S_Nf\|_2^2.
\]
给定 $\varepsilon>0$，推论~\ref{corollary:7-2-1-2} 给出三角多项式 $q$ 使
$\|f-q\|_2<\varepsilon$；当 $N$ 包含 $q$ 的全部频率时，上式说明
$\|f-S_Nf\|_2<\varepsilon$。故 $S_Nf\to f$ 于 $L^2$。在
\[
 \|S_Nf\|_2^2=\sum_{|n|\leq N}|\widehat f(n)|^2
\]
中令 $N\to\infty$，便得 Parseval 范数等式。对 $f+g$、$f+ig$ 使用范数等式，或直接把
$S_Nf,S_Ng$ 的有限和内积取极限，得到式~\eqref{eq:7-2-circle-parseval} 的内积版本。

尚需证明满射。任取 $a=(a_n)\in\ell^2(\Z)$，令
$p_N=\sum_{|n|\leq N}a_ne_n$。若 $M<N$，正交性给出
\[
 \|p_N-p_M\|_2^2=\sum_{M<|n|\leq N}|a_n|^2\longrightarrow0.
\]
所以 $(p_N)$ 在 $L^2(\T)$ 中为 Cauchy。$L^2$ 的完备性给出
$f\in L^2$ 使 $p_N\to f$。固定 $n$ 并取 $N\geq|n|$；Cauchy--Schwarz 不等式给出
\[
 |\widehat f(n)-a_n|
 =|\langle f-p_N,e_n\rangle|
 \leq\|f-p_N\|_2\longrightarrow0.
\]
因此 $\widehat f(n)=a_n$ 对每个 $n$ 成立。这一步同时使用了 $L^2$ 完备性，并证明 Fourier
映射的像不是 $\ell^2$ 的某个真闭子空间。
\end{proof}

\begin{proposition}[Fourier 唯一性]\label{proposition:7-2-1-4}
若 $\mu$ 是 $\T$ 上有限复 Borel 测度，且
$\widehat\mu(n)=0$ 对每个 $n\in\Z$ 成立，则 $\mu=0$。因此两个有限复测度具有相同的全部
三角矩，当且仅当它们相等。
\end{proposition}

\begin{proof}
若 $p(x)=\sum_{|n|\leq N}c_ne^{inx}$，则
\[
 \int_\T p\dd\mu=\sum_{|n|\leq N}c_n\widehat\mu(-n)=0.
\]
给定 $h\in C(\T)$，由推论~\ref{corollary:7-2-1-2} 可取三角多项式 $p_k$ 使
$\|p_k-h\|_\infty\to0$。全变差估计给出
\[
 \left|\int_\T h\dd\mu\right|
 \leq\left|\int_\T(h-p_k)\dd\mu\right|
 \leq\|h-p_k\|_\infty |\mu|(\T)\longrightarrow0.
\]
故 $\mu$ 消去所有连续函数。为从这里恢复测度本身，写极分解
$\dd\mu=u\dd|\mu|$，其中 $|u|=1$ a.e.。连续函数在
$L^1(|\mu|)$ 中稠密，故可取 $h_k\in C(\T)$ 使
$h_k\to\overline u$ 于 $L^1(|\mu|)$。于是
\[
 0=\lim_{k\to\infty}\int h_k\dd\mu
  =\int\overline u\,u\dd|\mu|=|\mu|(\T).
\]
所以 $|\mu|=0$，从而 $\mu=0$。对两个测度应用此结论于其差，即得最后一句。
\end{proof}

图~\ref{fig:7-2-dirichlet-fejer} 把两个核的差异画在同一尺度上。Dirichlet 核的正负旁瓣使
$L^1$ 质量累积；\Fejer{} 核只有非负主峰，峰外质量随 $N$ 消失。恢复机制不是“删除高频”，而是给
第 $n$ 个已观测矩乘上线性权重 $(1-|n|/(N+1))_+$。

\begin{figure}[htbp]
\centering
\begin{tikzpicture}[x=0.9cm,y=0.24cm]
  \begin{scope}
    \clip (-3.35,-7) rectangle (3.35,18.5);
    \draw[->,BookInk!70] (-3.3,0)--(3.3,0) node[right] {$x$};
    \draw[->,BookInk!70] (0,-6.5)--(0,18) node[above] {$D_8(x)$};
    \draw[BookWarning,thick,domain=-3.14:-0.04,samples=180]
      plot (\x,{sin(8.5*deg(\x))/sin(0.5*deg(\x))});
    \draw[BookWarning,thick,domain=0.04:3.14,samples=180]
      plot (\x,{sin(8.5*deg(\x))/sin(0.5*deg(\x))});
    \fill[BookWarning] (0,17) circle[radius=1.4pt];
    \node[book warning node,anchor=north west] at (-3.2,17.2)
      {符号振荡\\$\|D_N\|_1\asymp\log N$};
  \end{scope}
  \begin{scope}[xshift=8.2cm]
    \clip (-3.35,-1.3) rectangle (3.35,10.3);
    \draw[->,BookInk!70] (-3.3,0)--(3.3,0) node[right] {$x$};
    \draw[->,BookInk!70] (0,-1)--(0,10) node[above] {$F_8(x)$};
    \draw[BookAccent,thick,domain=-3.14:-0.04,samples=180]
      plot (\x,{(sin(4.5*deg(\x))/sin(0.5*deg(\x)))^2/9});
    \draw[BookAccent,thick,domain=0.04:3.14,samples=180]
      plot (\x,{(sin(4.5*deg(\x))/sin(0.5*deg(\x)))^2/9});
    \fill[BookAccent] (0,9) circle[radius=1.4pt];
    \node[book strong node,anchor=north west] at (-3.2,9.7)
      {正峰\\$\int F_N\dd m=1$};
  \end{scope}
  \node[book node] (mom) at (1.1,-10.8) {三角矩\\$\widehat f(n)$};
  \node[book warm node,right=1.4cm of mom] (weight)
    {\Fejer{} 权重\\$\left(1-|n|/(N+1)\right)_+$};
  \node[book strong node,right=1.4cm of weight] (rec) {恢复\\$\sigma_Nf\to f$};
  \draw[book arrow] (mom)--(weight);
  \draw[book accent arrow] (weight)--(rec);
\end{tikzpicture}
\caption{Dirichlet 振荡与 \Fejer{} 正平均。下方流程说明正核如何把三角矩变成稳定恢复。}
\label{fig:7-2-dirichlet-fejer}
\end{figure}

\begin{proposition}[周期热流与循环随机游走]\label{proposition:7-2-1-periodic-evolution}
\begin{enumerate}
  \item 对 $f\in L^2(\T)$ 和 $t\geq0$，由
  \[
    \widehat{H_tf}(n)=e^{-n^2t}\widehat f(n)
  \]
  定义的算子族 $(H_t)_{t\geq0}$ 是 $L^2(\T)$ 上的强连续收缩半群。对每个 $t>0$，
  $u(t,\cdot)=H_tf$ 有光滑代表元，并满足
  \[
    \partial_tu=\partial_x^2u,
    \qquad H_tf\longrightarrow f\quad\text{于 }L^2(\T)\text{ 中}\quad(t\downarrow0).
  \]
  \item 设 $M\geq1$、$G=\Z/M\Z$，$p$ 是 $G$ 上的概率分布，并令
  \[
    (Pf)(x)=\sum_{y\in G}p(y)f(x+y),
    \qquad \chi_k(x)=e^{2\pi ikx/M}.
  \]
  若
  \[
    \widehat f(k)=\frac1M\sum_{x\in G}f(x)\overline{\chi_k(x)},
    \qquad
    \lambda_k=\sum_{y\in G}p(y)\chi_k(y),
  \]
  则对每个整数 $m\geq0$，
  \begin{equation}\label{eq:7-2-cyclic-walk-diagonalization}
    P^mf=\sum_{k=0}^{M-1}\lambda_k^m\widehat f(k)\chi_k.
  \end{equation}
\end{enumerate}
\end{proposition}

\begin{proof}
Plancherel 等式给出
\[
 \norm{H_tf}_2^2
 =\sum_{n\in\Z}e^{-2n^2t}|\widehat f(n)|^2
 \leq\norm f_2^2.
\]
乘子恒等式
$e^{-n^2(s+t)}=e^{-n^2s}e^{-n^2t}$ 给出 $H_{s+t}=H_sH_t$。又由支配收敛，
\[
 \norm{H_tf-f}_2^2
 =\sum_{n\in\Z}|e^{-n^2t}-1|^2|\widehat f(n)|^2\longrightarrow0.
\]
半群律与收缩性把零时刻的连续性推广到每个 $t\geq0$。
若 $t>0$ 且 $r\geq0$，Cauchy--Schwarz 不等式给
\[
 \sum_{n\in\Z}|n|^re^{-n^2t}|\widehat f(n)|
 \leq\left(\sum_{n\in\Z}|n|^{2r}e^{-2n^2t}\right)^{1/2}\norm f_2<\infty.
\]
因此 Fourier 级数可任意次逐项对 $x$ 微分；在任意 $t\geq t_0>0$ 上还可逐项对 $t$
微分。乘子的导数为 $-n^2e^{-n^2t}$，故 $\partial_tH_tf=\partial_x^2H_tf$。

在有限群 $G$ 上，对 $0\leq k,\ell<M$，有限几何和给出
\[
 \frac1M\sum_{x\in G}\chi_k(x)\overline{\chi_\ell(x)}
 =\begin{cases}
 1,&k=\ell,\\
 \displaystyle\frac{1-e^{2\pi i(k-\ell)}}
 {M\bigl(1-e^{2\pi i(k-\ell)/M}\bigr)}=0,&k\ne\ell.
 \end{cases}
\]
因此这些特征标构成 $M$ 维空间 $\C^G$ 在归一化计数测度下的正交基，从而
$f=\sum_k\widehat f(k)\chi_k$。直接计算得
\[
 P\chi_k(x)=\sum_{y\in G}p(y)\chi_k(x+y)
 =\lambda_k\chi_k(x).
\]
对有限 Fourier 展开逐项应用 $P^m$，即得
\eqref{eq:7-2-cyclic-walk-diagonalization}。例如，当 $M\geq3$ 且
$p(1)=p(-1)=1/2$ 时，
\[
 \lambda_k=\cos\frac{2\pi k}{M}.
\]
若 $M$ 为奇数，则 $k\neq0$ 时 $|\lambda_k|<1$，故 $P^mf$ 趋于 $f$ 的平均值；若 $M$ 为偶数，则
$\lambda_{M/2}=-1$，对应模式显示出二周期振荡。
\end{proof}

\begin{exercise}
\begin{enumerate}
  \item 从 $S_Nf=f*D_N$ 与 $\sigma_Nf=(N+1)^{-1}\sum_{j=0}^NS_jf$ 出发，逐步推出
        \[
        D_N(x)=\frac{\sin((N+\tfrac12)x)}{\sin(x/2)},\qquad
        F_N(x)=\frac1{N+1}\left(\frac{\sin((N+1)x/2)}{\sin(x/2)}\right)^2,
        \]
        并在 $x=0$ 处用连续延拓解释公式。
  \item 对每个 $\delta\in(0,\pi)$ 给出
        $\sup_{\delta\le |x|\le\pi}F_N(x)$ 的显式上界，由此证明
        $\int_{|x|\ge\delta}F_N\dd m\to0$；同时直接核对 $F_N\ge0$ 与
        $\int_\T F_N\dd m=1$。
  \item 给定 $f\in C(\T)$ 与 $\varepsilon>0$，把
        $\sigma_Nf(x)-f(x)$ 的积分分成 $|t|<\delta$ 与 $|t|\ge\delta$ 两段，完整证明
        $\norm{\sigma_Nf-f}_\infty\to0$。再指出把一致连续性替换成 $L^p$ 平移连续性后，哪一步给出
        $1\le p<\infty$ 的 $L^p$ 收敛。
  \item 设有限复 Borel 测度 $\mu$ 的全部 Fourier 系数为零。先证明它消去每个三角多项式，再用
        \Fejer{} 稠密性与全变差估计证明它消去 $C(\T)$；最后写出极分解逼近步骤，推出
        $|\mu|(\T)=0$。
\end{enumerate}
\end{exercise}

圆周上的频率是整数，Fourier 映射把 $L^2(\T)$ 完整地离散化为 $\ell^2(\Z)$。在
$\R^d$ 上，频率变成连续变量；相应的求和、矩阵坐标和离散 Parseval 等式将分别变成积分、乘子和
Plancherel 等式。


\subsection{\texorpdfstring{$\R^d$}{R d} Fourier 变换、卷积与反演}\label{subsec:7-2-2}
\index{Fourier 变换}\index{Riemann--Lebesgue 定理}\index{Fourier 反演}

平移不变算子的共同特征是：指数函数把平移变成一个相位因子。若
$\tau_af(x)=f(x-a)$，则形式上
\[
 \int_{\R^d}f(x-a)e^{-2\pi ix\cdot\xi}\dd x
 =e^{-2\pi ia\cdot\xi}\int_{\R^d}f(x)e^{-2\pi ix\cdot\xi}\dd x.
\]
因此，卷积在频率侧应成为乘法，微分应成为多项式乘子。以下三个计算先展示这套规则的力量及其
边界；本小节始终使用指数 $e^{-2\pi ix\cdot\xi}$。

\begin{example}[Gaussian]\label{ex:7-2-2-1}
先在一维令 $g(x)=e^{-\pi x^2}$，并记
$G(\xi)=\int_\R g(x)e^{-2\pi ix\xi}\dd x$。对 $|h|\leq1$，差商恒等式给
\[
 \left|g(x)e^{-2\pi ix\xi}
       \frac{e^{-2\pi ixh}-1}{h}\right|
 \leq2\pi|x|e^{-\pi x^2},
\]
而右端可积，故支配收敛允许对 $\xi$ 求导。又因
$g'(x)=-2\pi xg(x)$，分部积分的边界项为零，故
\begin{align*}
 G'(\xi)
 &=-2\pi i\int_\R xg(x)e^{-2\pi ix\xi}\dd x\\
 &=i\int_\R g'(x)e^{-2\pi ix\xi}\dd x
 =-2\pi\xi G(\xi).
\end{align*}
Gaussian 积分给出 $G(0)=\int e^{-\pi x^2}\dd x=1$。把微分方程乘以
$e^{\pi\xi^2}$，可见
\[
 \frac{\dd}{\dd\xi}\bigl(e^{\pi\xi^2}G(\xi)\bigr)=0;
\]
代入初值得到
\[
 G(\xi)=e^{-\pi\xi^2}.
\]
对 $g_a(x)=e^{-\pi a|x|^2}$，$a>0$，逐坐标使用 Fubini，再作伸缩
$y=\sqrt a\,x$，得到
\begin{equation}\label{eq:7-2-gaussian-transform}
 \widehat{g_a}(\xi)=a^{-d/2}e^{-\pi|\xi|^2/a}.
\end{equation}
Gaussian 因而在空间侧与频率侧具有相同形状，宽度则互为倒数。
\end{example}

\begin{example}[区间指标]\label{ex:7-2-2-2}
在一维令 $I=[-1/2,1/2]$。当 $\xi\neq0$ 时
\[
 \widehat{\1_I}(\xi)
 =\int_{-1/2}^{1/2}e^{-2\pi ix\xi}\dd x
 =\frac{e^{-\pi i\xi}-e^{\pi i\xi}}{-2\pi i\xi}
 =\frac{\sin(\pi\xi)}{\pi\xi};
\]
在 $\xi=0$ 处连续延拓值为 $1$。该 sinc 函数不属于 $L^1(\R)$。事实上，对每个整数
$k\geq1$，在 $[k+1/6,k+5/6]$ 上有 $|\sin(\pi\xi)|\geq1/2$，从而
\[
 \int_1^\infty\left|\frac{\sin(\pi\xi)}{\pi\xi}\right|\dd\xi
 \geq\frac1{2\pi}\sum_{k=1}^\infty
   \frac{2/3}{k+5/6}=\infty.
\]
所以 $f\in L^1$ 并不保证 $\widehat f\in L^1$。区间边界的跳跃只产生 $|\xi|^{-1}$ 级衰减，
未经正则化的反演积分不绝对收敛。
\end{example}

\begin{example}[卷积平滑]\label{ex:7-2-2-3}
仍令 $I=[-1/2,1/2]$。对每个 $x$，两个区间 $I$ 与 $x-I$ 的交长给出
\[
 (\1_I*\1_I)(x)
 =\int_\R\1_I(x-y)\1_I(y)\dd y
 =(1-|x|)_+.
\]
这次卷积把两个跳跃函数变成连续的帐篷函数。另一方面，稍后证明的卷积公式给出
\[
 \widehat{(1-|\cdot|)_+}(\xi)
 =\left(\frac{\sin(\pi\xi)}{\pi\xi}\right)^2.
\]
空间侧多一次卷积，频率侧就多一个 sinc 衰减因子；平方后的函数已经属于 $L^1(\R)$。
\end{example}

\begin{definition}[Fourier 变换约定]\label{def:7-2-2-1}
对 $f\in L^1(\R^d)$，定义
\[
 \widehat f(\xi)=\mathcal Ff(\xi)
 =\int_{\R^d}f(x)e^{-2\pi ix\cdot\xi}\dd x.
\]
对 $g\in L^1(\R^d)$，定义反 Fourier 积分
\[
 \check g(x)=\mathcal F^{-1}g(x)
 =\int_{\R^d}g(\xi)e^{2\pi ix\cdot\xi}\dd\xi.
\]
对有限复 Borel 测度 $\mu$，相同约定给出
$\widehat\mu(\xi)=\int e^{-2\pi ix\cdot\xi}\dd\mu(x)$。
\end{definition}

文献中常把 $2\pi$ 放在逆变换的常数里。三种常用写法的换算如下；在本书余下部分只使用第一行。
\begin{center}
\renewcommand{\arraystretch}{1.25}
\small
\begin{tabularx}{\textwidth}{@{}p{2.0cm}XX@{}}
\toprule
约定 & 正变换 & 逆变换\\
\midrule
本书 & $\widehat f(\xi)=\int f(x)e^{-2\pi ix\cdot\xi}\dd x$
     & $f(x)=\int\widehat f(\xi)e^{2\pi ix\cdot\xi}\dd\xi$\\
角频率 & $F_0(\omega)=\int f(x)e^{-ix\cdot\omega}\dd x$
     & $f(x)=(2\pi)^{-d}\int F_0(\omega)e^{ix\cdot\omega}\dd\omega$\\
酉角频率 & $F_u(\omega)=(2\pi)^{-d/2}\int f(x)e^{-ix\cdot\omega}\dd x$
     & $f(x)=(2\pi)^{-d/2}\int F_u(\omega)e^{ix\cdot\omega}\dd\omega$\\
\bottomrule
\end{tabularx}
\end{center}
其中 $F_0(\omega)=\widehat f(\omega/(2\pi))$，且
$F_u=(2\pi)^{-d/2}F_0$。常数的所有差异都来自换元 $\omega=2\pi\xi$。

\begin{definition}[调制、平移与伸缩]\label{def:7-2-2-2}
对 $a,b\in\R^d$ 及可逆线性映射 $A:\R^d\to\R^d$，定义
\[
 \tau_af(x)=f(x-a),\qquad
 M_bf(x)=e^{2\pi ib\cdot x}f(x),\qquad
 D_Af(x)=f(Ax).
\]
\end{definition}

\begin{definition}[可反演 $L^1$ 类]\label{def:7-2-2-3}
称 $f$ 属于可反演 $L^1$ 类，如果
\[
 f\in L^1(\R^d),\qquad \widehat f\in L^1(\R^d).
\]
\end{definition}

\begin{theorem}[Riemann--Lebesgue]\label{theorem:7-2-2-1}
若 $f\in L^1(\R^d)$，则 $\widehat f\in C_0(\R^d)$，且
\[
 \|\widehat f\|_\infty\leq\|f\|_1.
\]
\end{theorem}

\begin{proof}
三角不等式立即给出
$|\widehat f(\xi)|\leq\int|f|=\|f\|_1$。若 $\xi_k\to\xi$，则积分中的指数逐点收敛，且
\[
 |f(x)e^{-2\pi ix\cdot\xi_k}|\leq|f(x)|.
\]
支配收敛定理给出 $\widehat f(\xi_k)\to\widehat f(\xi)$，故 $\widehat f$ 连续。

先设 $g\in C_c(\R^d)$。给 $\xi\neq0$，选坐标 $j$ 使
$|\xi_j|\geq|\xi|/\sqrt d$，并令 $h=e_j/(2\xi_j)$。于是
$e^{-2\pi ih\cdot\xi}=-1$。在线性换元后，
\begin{align*}
 2\widehat g(\xi)
 &=\int_{\R^d}\bigl(g(x)-g(x-h)\bigr)e^{-2\pi ix\cdot\xi}\dd x,\\
 2|\widehat g(\xi)|&\leq\|g-\tau_hg\|_1.
\end{align*}
当 $|\xi|\to\infty$ 时，$|h|\leq\sqrt d/(2|\xi|)\to0$；$L^1$ 平移连续性使右端趋于零。

对一般 $f\in L^1$，给定 $\varepsilon>0$，取 $g\in C_c$ 使
$\|f-g\|_1<\varepsilon$。由第一步的上界，
\[
 |\widehat f(\xi)|
 \leq\|f-g\|_1+|\widehat g(\xi)|<2\varepsilon
\]
对所有充分大的 $|\xi|$ 成立。因此 $\widehat f$ 在无穷远消失。
\end{proof}

\begin{proposition}[代数公式]\label{proposition:7-2-2-2}
在下列各项所写的可积条件下，Fourier 变换满足：
\begin{enumerate}
  \item 若 $f,g\in L^1$，则
  $\widehat{f*g}=\widehat f\,\widehat g$；
  \item 若 $f\in W^{1,1}(\R^d)$，则
  $\widehat{\partial_jf}(\xi)=2\pi i\xi_j\widehat f(\xi)$；
  \item 若 $f,x_jf\in L^1$，则 $\widehat f$ 关于 $\xi_j$ 连续可微，且
  \[
    \widehat{x_jf}(\xi)=-\frac1{2\pi i}\partial_{\xi_j}\widehat f(\xi);
  \]
  \item 若 $f\in L^1(\R^d)$，则 $\tau_af,M_bf,D_Af\in L^1(\R^d)$，并且
  \[
    \widehat{\tau_af}(\xi)=e^{-2\pi ia\cdot\xi}\widehat f(\xi),
    \qquad
    \widehat{M_bf}(\xi)=\widehat f(\xi-b),
  \]
  \[
    \widehat{D_Af}(\xi)=|\det A|^{-1}\widehat f(A^{-T}\xi).
  \]
\end{enumerate}
\end{proposition}

\begin{proof}
若 $f,g\in L^1$，则
\[
 \int\!\!\int |f(x-y)g(y)|\dd y\dd x=\|f\|_1\|g\|_1<\infty.
\]
所以 Fubini 定理与换元 $z=x-y$ 给出
\begin{align*}
 \widehat{f*g}(\xi)
 &=\int\!\!\int f(x-y)g(y)e^{-2\pi ix\cdot\xi}\dd y\dd x\\
 &=\int\!\!\int f(z)g(y)e^{-2\pi i(z+y)\cdot\xi}\dd z\dd y
 =\widehat f(\xi)\widehat g(\xi).
\end{align*}

若 $f\in W^{1,1}$，取 $\chi\in C_c^\infty$ 使
$0\leq\chi\leq1$、$\chi=1$ 于单位球，并令 $\chi_R(x)=\chi(x/R)$。弱导数的积分分部公式给出
\begin{align*}
 \int \partial_jf(x)e^{-2\pi ix\cdot\xi}\chi_R(x)\dd x
 &=2\pi i\xi_j\int f(x)e^{-2\pi ix\cdot\xi}\chi_R(x)\dd x\\
 &\quad-\int f(x)e^{-2\pi ix\cdot\xi}\partial_j\chi_R(x)\dd x.
\end{align*}
前两项由支配收敛分别趋于 $\widehat{\partial_jf}(\xi)$ 与
$2\pi i\xi_j\widehat f(\xi)$；最后一项绝对值不超过
\[
 R^{-1}\|\partial_j\chi\|_\infty\|f\|_1,
\]
故趋于零。这证明微分公式。

若 $x_jf\in L^1$，差商恒等式以及
$|(e^{-it}-1)/t|\leq1$ 给出一个由 $2\pi|x_jf(x)|$ 控制的支配函数。因此
\[
 \partial_{\xi_j}\widehat f(\xi)
 =-2\pi i\int x_jf(x)e^{-2\pi ix\cdot\xi}\dd x
 =-2\pi i\widehat{x_jf}(\xi).
\]
同一支配收敛论证说明导数连续。

最后，平移、调制与线性伸缩的 $L^1$ 归属分别由平移不变性、指数因子的模为 $1$ 以及线性换元
得到。前两式只需分别作 $y=x-a$ 和合并指数。在线性伸缩中令 $y=Ax$，便有
\begin{align*}
 \widehat{D_Af}(\xi)
 &=|\det A|^{-1}\int f(y)e^{-2\pi i(A^{-1}y)\cdot\xi}\dd y\\
 &=|\det A|^{-1}\widehat f(A^{-T}\xi).
\end{align*}
这里 $|\det A|^{-1}$ 来自换元公式中的雅可比行列式因子，而 $A^{-T}$ 由
$(A^{-1}y)\cdot\xi=y\cdot A^{-T}\xi$ 决定。
\end{proof}

\begin{theorem}[Fourier 反演]\label{theorem:7-2-2-3}
若 $f,\widehat f\in L^1(\R^d)$，则在所选代表 $f$ 的每个 Lebesgue 点 $x$，
\begin{equation}\label{eq:7-2-fourier-inversion}
 f(x)=\int_{\R^d}\widehat f(\xi)e^{2\pi ix\cdot\xi}\dd\xi.
\end{equation}
特别地，等式在每个连续点成立，而且 $f$ 由 $\widehat f$ 唯一决定。
\end{theorem}

\begin{proof}
对 $\varepsilon>0$，令
\[
 \phi_\varepsilon(x)=\varepsilon^{-d/2}e^{-\pi|x|^2/\varepsilon}.
\]
由例~\ref{ex:7-2-2-1}，
$\widehat{\phi_\varepsilon}(\xi)=e^{-\pi\varepsilon|\xi|^2}$，且
$\int\phi_\varepsilon=1$。由于
\[
 \int\!\!\int |f(y)|e^{-\pi\varepsilon|\xi|^2}\dd\xi\dd y
 =\varepsilon^{-d/2}\|f\|_1<\infty,
\]
Fubini 定理允许交换积分，得到 Gaussian 正则化恒等式
\begin{align}
 I_\varepsilon(x)
 &:=\int\widehat f(\xi)e^{2\pi ix\cdot\xi}
                 e^{-\pi\varepsilon|\xi|^2}\dd\xi\notag\\
 &=\int f(y)\left(\int e^{2\pi i(x-y)\cdot\xi}
                 e^{-\pi\varepsilon|\xi|^2}\dd\xi\right)\dd y\notag\\
 &=(f*\phi_\varepsilon)(x).
 \label{eq:7-2-gaussian-regularization}
\end{align}
写 $\delta=\sqrt\varepsilon$，便有
$\phi_\varepsilon(x)=\delta^{-d}e^{-\pi|x/\delta|^2}$。所以它正是质量一的径向递减伸缩核；
第四章的 Lebesgue 点近似恒等定理
\ref{theorem:4-4-2-2} 因而给出
$(f*\phi_\varepsilon)(x)\to f(x)$，只要 $x$ 是 $f$ 的 Lebesgue 点。
另一方面，$\widehat f\in L^1$，且 Gaussian 乘子绝对值不超过 $1$ 并逐点趋于 $1$；支配收敛给出
\[
 I_\varepsilon(x)\longrightarrow
 \int\widehat f(\xi)e^{2\pi ix\cdot\xi}\dd\xi.
\]
在式~\eqref{eq:7-2-gaussian-regularization} 两边取极限，即得反演公式。

右端是连续函数，因为对 $x_k\to x$ 可用 $|\widehat f|$ 作支配。Lebesgue 微分定理说明几乎每点是
Lebesgue 点，故这个连续函数与 $f$ a.e. 相等。两个连续函数若 a.e. 相等，则其差的非零点集既开又为
零测，只能为空；所以连续代表唯一。若两个可反演函数的 Fourier 变换相同，反演公式给出其连续代表
相同，因而两个 $L^1$ 等价类相同。
\end{proof}

\begin{proposition}[测度唯一性]\label{proposition:7-2-2-4}
若 $\mu,\nu$ 是 $\R^d$ 上有限复 Borel 测度，且
$\widehat\mu(\xi)=\widehat\nu(\xi)$ 对每个 $\xi\in\R^d$ 成立，则 $\mu=\nu$。
\end{proposition}

\begin{proof}
令 $\lambda=\mu-\nu$，则 $\widehat\lambda=0$。用反演证明中的 Gaussian 定义
\[
 h_\varepsilon(x)=\int_{\R^d}\phi_\varepsilon(x-y)\dd\lambda(y).
\]
Tonelli 定理给出
$\|h_\varepsilon\|_1\leq|\lambda|(\R^d)$。Gaussian 的一致连续性与全变差估计说明
$h_\varepsilon$ 连续。再用 Fubini 及命题~\ref{proposition:7-2-2-2} 的同一计算，
\[
 \widehat{h_\varepsilon}(\xi)
 =e^{-\pi\varepsilon|\xi|^2}\widehat\lambda(\xi)=0.
\]
于是 $h_\varepsilon,\widehat{h_\varepsilon}\in L^1$；定理~\ref{theorem:7-2-2-3} 作用于连续函数
$h_\varepsilon$，给出 $h_\varepsilon(x)=0$ 对每个 $x$ 成立。

任取 $\psi\in C_c(\R^d)$。Fubini 定理和 $\phi_\varepsilon(-z)=\phi_\varepsilon(z)$ 给出
\[
 0=\int\psi(x)h_\varepsilon(x)\dd x
 =\int(\psi*\phi_\varepsilon)(y)\dd\lambda(y).
\]
Gaussian 是近似恒等族，而 $\psi\in C_0$，故
$\psi*\phi_\varepsilon\to\psi$ 一致。由
\[
 \left|\int(\psi*\phi_\varepsilon-\psi)\dd\lambda\right|
 \leq\|\psi*\phi_\varepsilon-\psi\|_\infty|\lambda|(\R^d)
\]
可知 $\int\psi\dd\lambda=0$。最后写
$\dd\lambda=u\dd|\lambda|$，并用 $C_c(\R^d)$ 在
$L^1(|\lambda|)$ 中的稠密性逼近 $\overline u$，得到
$|\lambda|(\R^d)=0$。所以 $\lambda=0$，即 $\mu=\nu$。
\end{proof}

图~\ref{fig:7-2-convolution-inversion} 同时记录两条路径。上方交换方块把空间卷积送到频率乘法；
下方则说明反演并非把振荡积分直接“算回去”，而是先乘 Gaussian、回到空间侧的近似恒等，再令
$\varepsilon\downarrow0$。积分交换在固定 $\varepsilon>0$ 时绝对可积，这正是正则化的作用。

\begin{figure}[htbp]
\centering
\begin{tikzpicture}[node distance=1.25cm and 2.0cm]
  \node[book node] (fg) {$f,\ g\in L^1$};
  \node[book warm node,right=of fg] (conv) {空间卷积\\$f*g$};
  \node[book node,below=of fg] (hats) {$\widehat f,\ \widehat g$};
  \node[book strong node,right=of hats] (prod) {频率乘法\\$\widehat f\,\widehat g$};
  \draw[book arrow] (fg)--node[above,book label]{卷积} (conv);
  \draw[book arrow] (fg)--node[left,book label]{Fourier} (hats);
  \draw[book arrow] (conv)--node[right,book label]{Fourier} (prod);
  \draw[book accent arrow] (hats)--node[below,book label]{逐点乘} (prod);

  \node[book node,below=1.8cm of hats] (obs) {观测\\$\widehat f$};
  \node[book warm node,right=1.8cm of obs] (reg)
    {Gaussian 正则化\\$e^{-\pi\varepsilon|\xi|^2}\widehat f$};
  \node[book strong node,below=1.5cm of reg] (ai)
    {空间近似恒等\\$f*\phi_\varepsilon$};
  \node[book strong node,right=1.8cm of ai] (limit) {Lebesgue 点\\$f(x)$};
  \draw[book arrow] (obs)--(reg);
  \draw[book accent arrow] (reg)--node[right,book label,align=center]{绝对可积\\反变换} (ai);
  \draw[book accent arrow] (ai)--node[above,book label]{$\varepsilon\downarrow0$} (limit);
\end{tikzpicture}
\caption{卷积--乘法交换方块与 Gaussian 正则化反演路径。}
\label{fig:7-2-convolution-inversion}
\end{figure}

这项测度唯一性是第五章“特征函数决定概率分布”的有限复测度版本。它只回答相同变换能否来自两个
不同对象；若一列概率变换逐点收敛，要进一步推出概率弱收敛，仍需紧性以及极限在零点连续，
不能从唯一性单独推出存在性。

\begin{exercise}
\begin{enumerate}
  \item 从一维函数 $g(x)=e^{-\pi x^2}$ 的微分方程证明 $\widehat g=g$，再用 Fubini 与伸缩换元
        推出
        $\widehat{e^{-\pi a|x|^2}}(\xi)=a^{-d/2}e^{-\pi|\xi|^2/a}$；逐项注明支配函数与
        雅可比行列式因子。
  \item 在本章的 $2\pi$ 约定下，对 $f\in L^1(\R^d)$ 从定义证明平移、调制与可逆线性伸缩公式：
        \[
        \begin{aligned}
        \widehat{f(\,\cdot-a)}(\xi)
          &=e^{-2\pi ia\cdot\xi}\widehat f(\xi),\\
        \widehat{e^{2\pi i\eta\cdot x}f(x)}(\xi)
          &=\widehat f(\xi-\eta),\\
        \widehat{f(A\,\cdot)}(\xi)
          &=|\det A|^{-1}\widehat f(A^{-T}\xi).
        \end{aligned}
        \]
  \item 对 $f,g\in L^1(\R^d)$ 从头验证定义 $f*g$ 所需的 a.e. 可积性，并写出允许交换积分的
        绝对可积上界，由此证明 $\widehat{f*g}=\widehat f\,\widehat g$。
  \item 令 $\phi_\varepsilon(x)=\varepsilon^{-d/2}e^{-\pi|x|^2/\varepsilon}$。在
        $f,\widehat f\in L^1$ 的假设下，完整证明
        \[
        \int\widehat f(\xi)e^{2\pi ix\cdot\xi}e^{-\pi\varepsilon|\xi|^2}\dd\xi
        =(f*\phi_\varepsilon)(x),
        \]
        并分别说明 Lebesgue 点近似恒等定理与支配收敛在令 $\varepsilon\downarrow0$ 时负责哪一边。
\end{enumerate}
\end{exercise}


\subsection{Schwartz 空间、Plancherel 与热/Poisson 半群}\label{subsec:7-2-3}
\index{Schwartz 空间}\index{Fourier 乘子}\index{热半群}\index{Poisson 半群}

$L^1$ 反演的附加条件 $\widehat f\in L^1$ 并不自动成立。为了得到一个对微分、乘多项式和
Fourier 变换均封闭的测试函数空间，需要同时控制无穷远处的衰减与所有阶导数。Schwartz 空间用一族
半范数同时表达这些要求，而不将它们归结为单个范数。随后，$L^2$ 变换将由这个稠密测试函数空间完备化得到；
这与圆周上从三角多项式完备化到 $L^2(\T)$ 是同一结构。

若 $\alpha=(\alpha_1,\ldots,\alpha_d)\in(\Z_{\geq0})^d$，沿用
$|\alpha|=\sum_j\alpha_j$、$x^\alpha=\prod_jx_j^{\alpha_j}$ 与
$\partial^\alpha=\prod_j\partial_j^{\alpha_j}$ 的多重指标记号。

\begin{definition}[Schwartz 空间]\label{def:7-2-3-1}
Schwartz 空间 $\mathcal S(\R^d)$ 由所有 $f\in C^\infty(\R^d)$ 组成，并要求对每对多重指标
$\alpha,\beta$，
\[
 p_{\alpha,\beta}(f)
 :=\sup_{x\in\R^d}|x^\alpha\partial^\beta f(x)|<\infty.
\]
取遍所有 $p_{\alpha,\beta}$ 所生成的局部凸拓扑；亦即
$f_n\to f$ 当且仅当每个 $p_{\alpha,\beta}(f_n-f)\to0$。这族可数半范数给出
$\mathcal S$ 的 \Frechet{} 拓扑。
\end{definition}

为说明最后一句，枚举半范数为 $(p_k)_{k\geq1}$，则
\[
 d_{\mathcal S}(f,g)=\sum_{k=1}^\infty2^{-k}
   \min\{1,p_k(f-g)\}
\]
诱导上述拓扑。若 $(f_n)$ 关于该度量 Cauchy，则固定一个半范数 $p_k$ 和
$0<\varepsilon<1$ 后，当
$d_{\mathcal S}(f_n,f_m)<2^{-k}\varepsilon$ 时必有
$p_k(f_n-f_m)<\varepsilon$。因此每个 $(\partial^\beta f_n)$ 关于全空间上确界范数 Cauchy，存在
有界连续函数 $g_\beta$ 使
\[
 \|\partial^\beta f_n-g_\beta\|_\infty\longrightarrow0.
\]
置 $f=g_0$。对每个多重指标 $\gamma$、坐标 $j$ 与实数 $t$，逐坐标使用
\[
 \partial^\gamma f_n(x+te_j)-\partial^\gamma f_n(x)
 =\int_0^t\partial^{\gamma+e_j}f_n(x+se_j)\dd s.
\]
被积函数在连接 $x$ 与 $x+te_j$ 的线段上一致收敛；令 $n\to\infty$ 得
\[
 g_\gamma(x+te_j)-g_\gamma(x)
 =\int_0^tg_{\gamma+e_j}(x+se_j)\dd s.
\]
微积分基本定理于是给 $\partial_jg_\gamma=g_{\gamma+e_j}$。对 $|\gamma|$ 归纳可知
$f\in C^\infty$ 且 $\partial^\beta f=g_\beta$。

最后，$(x^\alpha\partial^\beta f_n)$ 在上确界范数中也是 Cauchy，故一致收敛到某个
$h_{\alpha,\beta}$；逐点取极限给
$h_{\alpha,\beta}=x^\alpha\partial^\beta f$。因此
\[
 p_{\alpha,\beta}(f_n-f)
 =\|x^\alpha\partial^\beta f_n-h_{\alpha,\beta}\|_\infty
 \longrightarrow0.
\]
所以极限仍在 $\mathcal S$ 中，且该度量完备。

Schwartz 函数属于每个 $L^p$。例如
\begin{equation}\label{eq:7-2-schwartz-l1-bound}
 \|h\|_1
 \leq C_d\sup_x(1+|x|)^{d+1}|h(x)|
 \leq C'_d\sum_{|\gamma|\leq d+1}p_{\gamma,0}(h).
\end{equation}
微分或乘以多项式只会移动半范数的指标；Leibniz 公式还说明两个 Schwartz 函数的乘积仍在
$\mathcal S$。卷积也保持该空间。事实上，对 $f,g\in\mathcal S$，微分可移入卷积，而
$x=(x-y)+y$ 给出
\[
 x^\alpha\partial^\beta(f*g)(x)
 =\sum_{\gamma\leq\alpha}\binom{\alpha}{\gamma}
 \bigl[(z^{\alpha-\gamma}\partial^\beta f(z))
       *(z^\gamma g(z))\bigr](x).
\]
每一项都是一个 $L^\infty$ 函数与一个 $L^1$ 函数的卷积，故一致有界；这证明
$p_{\alpha,\beta}(f*g)<\infty$。

\begin{definition}[缓增分布]\label{def:7-2-3-2}
Schwartz 空间上的连续线性泛函称为\emph{缓增分布}（tempered distribution），全体记作
$\mathcal S'(\R^d)$。
\end{definition}

\begin{proposition}[缓增分布的基本例子]\label{proposition:7-2-3-1}
下列对象都以自然方式给出 $\mathcal S'(\R^d)$ 中的元素。
\begin{enumerate}
  \item 若 $f\in L^1_{\mathrm{loc}}(\R^d)$，并且对某些 $C>0$ 与整数 $N\geq0$ 有
  $|f(x)|\leq C(1+|x|)^N$ a.e.，则
  \[
    \langle T_f,\varphi\rangle=\int_{\R^d}f(x)\varphi(x)\dd x.
  \]
  \item 每个 $f\in L^p(\R^d)$（$1\leq p\leq\infty$）按同一公式定义缓增分布。
  \item 每个多项式按第一项的公式定义缓增分布。
  \item 若复 Radon 测度 $\mu$ 对某些 $C,M>0$ 满足
  \[
    |\mu|(B(0,R))\leq C(1+R)^M\qquad(R\geq1),
  \]
  则 $\langle T_\mu,\varphi\rangle=\int\varphi\dd\mu$ 定义缓增分布。
\end{enumerate}
\end{proposition}

\begin{proof}
在第一种情形，取整数 $K>N+d$。则
\[
 |\langle T_f,\varphi\rangle|
 \leq C\sup_x(1+|x|)^K|\varphi(x)|
       \int_{\R^d}(1+|x|)^{N-K}\dd x.
\]
加权上确界由有限个 $p_{\alpha,0}$ 控制，所以 $T_f$ 连续；多项式是这一情形的特例。

若 $f\in L^p$，\Holder{} 不等式给出
\[
 |\langle T_f,\varphi\rangle|
 \leq \|f\|_p\|\varphi\|_{p'},
\]
而当 $1\leq q<\infty$ 时，
\[
 \|\varphi\|_q
 \leq\left(\int_{\R^d}(1+|x|)^{-q(d+1)}\dd x\right)^{1/q}
       \sup_x(1+|x|)^{d+1}|\varphi(x)|;
\]
$q=\infty$ 时则有 $\|\varphi\|_\infty=p_{0,0}(\varphi)$。因此
$\mathcal S$ 连续嵌入每个 $L^q$，所以上述泛函连续。

最后考虑测度。把空间分成 $B(0,2)$ 与环带
$2^k\leq|x|<2^{k+1}$。取整数 $K>M+1$，则
\begin{align*}
 \int|\varphi|\dd|\mu|
 &\leq C_{M,K}\sup_x(1+|x|)^K|\varphi(x)|
   \left(1+\sum_{k=1}^\infty2^{-k(K-M)}\right).
\end{align*}
右端仍由有限个 Schwartz 半范数控制，因而 $T_\mu$ 连续。
\end{proof}

\begin{definition}[Fourier 乘子]\label{def:7-2-3-3}
给定频率函数 $m(\xi)$，在使右端有意义的函数上定义形式算子
\[
 T_mf=\mathcal F^{-1}(m\widehat f).
\]
称 $m$ 为该算子的 Fourier 乘子。
\end{definition}

本节的两个基本乘子族是
\[
 m_t^{\mathrm{heat}}(\xi)=e^{-4\pi^2t|\xi|^2},
 \qquad
 m_t^{\mathrm{Pois}}(\xi)=e^{-2\pi t|\xi|},
 \qquad t>0.
\]
指数恒等式 $m_s m_t=m_{s+t}$ 预示相应算子族具有半群律；要把形式计算变成定理，还须识别其空间核
并证明在所用函数空间中连续。

\begin{theorem}[Schwartz 自同构]\label{theorem:7-2-3-1}
Fourier 变换是 $\mathcal S(\R^d)$ 上的连续线性双射，逆映射为反 Fourier 变换，而且逆映射也连续。
换言之，$\mathcal F:\mathcal S\to\mathcal S$ 是 \Frechet{} 空间自同构。
\end{theorem}

\begin{proof}
对 $f\in\mathcal S$，命题~\ref{proposition:7-2-2-2} 的微分与乘法公式可以任意次迭代，得到
\[
 \partial_\xi^\beta\widehat f
 =(-2\pi i)^{|\beta|}\widehat{x^\beta f},
\]
以及
\begin{equation}\label{eq:7-2-schwartz-seminorm-identity}
 \xi^\alpha\partial_\xi^\beta\widehat f(\xi)
 =(2\pi i)^{-|\alpha|}(-2\pi i)^{|\beta|}
   \widehat{\partial^\alpha(x^\beta f)}(\xi).
\end{equation}
取上确界，使用 $\|\widehat h\|_\infty\leq\|h\|_1$ 与
式~\eqref{eq:7-2-schwartz-l1-bound}，再展开 Leibniz 公式，得到
\begin{equation}\label{eq:7-2-schwartz-continuity}
 p_{\alpha,\beta}(\widehat f)
 \leq C_{\alpha,\beta,d}
 \sum_{(\gamma,\delta)\in E_{\alpha,\beta}}
 p_{\gamma,\delta}(f),
\end{equation}
其中 $E_{\alpha,\beta}$ 是一个有限指标集。因此 $\widehat f\in\mathcal S$，而且每个输出半范数由
有限个输入半范数控制；这正是 $\mathcal F$ 的连续性。

令 $Rg(x)=g(-x)$。由定义 $\mathcal Gg:=\check g=R\widehat g$，而
$p_{\alpha,\beta}(Rg)=p_{\alpha,\beta}(g)$；所以刚证的估计说明反 Fourier 算子
$\mathcal G=R\mathcal F$ 也连续地把 $\mathcal S$ 映入自身。又因
$f,\widehat f\in L^1$ 且 $f$ 处处连续，反演定理~\ref{theorem:7-2-2-3} 给出
\[
 \mathcal G\mathcal Ff=f\qquad(f\in\mathcal S).
\]
由 $\mathcal G=R\mathcal F$，故
$\mathcal G\mathcal F=I$ 等价于 $R\mathcal F^2=I$，即
$\mathcal F^2=R$。线性换元又给出 $\mathcal F R=R\mathcal F$。因此
\[
 \mathcal F\mathcal G
 =\mathcal F R\mathcal F=R\mathcal F^2=R^2=I.
\]
所以 $\mathcal F$ 与 $\mathcal G$ 互为连续逆映射。
\end{proof}

\begin{definition}[缓增分布的 Fourier 变换]\label{def:7-2-3-4}
对 $T\in\mathcal S'(\R^d)$，定义其 Fourier 变换 $\widehat T$ 为
\[
 \langle\widehat T,\varphi\rangle
 :=\langle T,\widehat\varphi\rangle
 \qquad(\varphi\in\mathcal S(\R^d)).
\]
\end{definition}

Schwartz 自同构定理保证 $\varphi\mapsto\widehat\varphi$ 连续，因此上式确实定义
$\mathcal S'$ 中的元素。本章只用这一转置定义计算 Dirac 测度和常数函数的 Fourier 变换；一般的
分布运算留待后续课程。

\begin{theorem}[Plancherel 定理]\label{theorem:7-2-3-2}
对 $f,g\in\mathcal S(\R^d)$，
\begin{equation}\label{eq:7-2-euclidean-parseval}
 \langle\widehat f,\widehat g\rangle_{L^2}
 =\langle f,g\rangle_{L^2},
 \qquad \|\widehat f\|_2=\|f\|_2.
\end{equation}
Fourier 变换唯一延拓为 $L^2(\R^d)$ 上的满射线性等距算子；其逆是反 Fourier 变换的
$L^2$ 延拓。因此该延拓是酉算子。
\end{theorem}

\begin{proof}
若 $f,g\in\mathcal S$，则 $f,g,\widehat f,\widehat g$ 都属于 $L^1\cap L^2$。特别地，
\[
 \int\!\!\int|\widehat f(\xi)g(x)|\dd x\dd\xi
 =\|\widehat f\|_1\|g\|_1<\infty.
\]
故 Fubini 定理与反演公式给出
\begin{align*}
 \langle\widehat f,\widehat g\rangle
 &=\int\widehat f(\xi)
   \left(\int\overline{g(x)}e^{2\pi ix\cdot\xi}\dd x\right)\dd\xi\\
 &=\int\overline{g(x)}
   \left(\int\widehat f(\xi)e^{2\pi ix\cdot\xi}\dd\xi\right)\dd x\\
 &=\int f(x)\overline{g(x)}\dd x
 =\langle f,g\rangle.
\end{align*}
取 $g=f$ 得等距式。

$C_c^\infty(\R^d)\subset\mathcal S$，而第四章的光滑紧支撑稠密性定理
\ref{theorem:4-3-1-3} 说明 $\mathcal S$ 在 $L^2$ 中稠密。给 $f\in L^2$，取
$f_n\in\mathcal S$ 使 $f_n\to f$ 于 $L^2$。等距式给出
\[
 \|\widehat f_n-\widehat f_m\|_2=\|f_n-f_m\|_2,
\]
所以 $(\widehat f_n)$ 在完备的 $L^2$ 中有极限。定义
$\mathcal F_2f=\lim_n\widehat f_n$。同一等式说明定义与逼近列无关、保持线性，并且
$\|\mathcal F_2f\|_2=\|f\|_2$；稠密性还给出这种连续延拓的唯一性。

反 Fourier 变换在 $\mathcal S$ 上也是等距，故以相同步骤延拓为
$\mathcal G_2:L^2\to L^2$。在稠密子空间 $\mathcal S$ 上已有
$\mathcal G_2\mathcal F_2=\mathcal F_2\mathcal G_2=I$。两边都是有界算子；对任意 $L^2$
函数取 Schwartz 逼近并令极限通过算子，得到这些恒等式在整个 $L^2$ 上成立。因此
$\mathcal F_2$ 满射，逆为 $\mathcal G_2$。
\end{proof}

定理中的 $L^2$ Fourier 变换一般不是绝对收敛积分。例如当
$d/2<\alpha\leq d$ 时，$(1+|x|)^{-\alpha}\in L^2(\R^d)\setminus L^1(\R^d)$。
对此类函数，$\mathcal F_2f$ 的含义是任意 Schwartz 逼近列的变换在 $L^2$ 中的极限；换用另一逼近列
不会改变所得等价类，但定理不承诺每个频率点都有绝对积分值。

\begin{proposition}[热半群]\label{proposition:7-2-3-3}
令
\[
 p_t(x)=(4\pi t)^{-d/2}e^{-|x|^2/(4t)},
 \qquad P_tf=p_t*f,\qquad t>0,
\]
并令 $P_0=I$。对每个固定的 $1\leq p<\infty$，$(P_t)_{t\geq0}$ 是
$L^p(\R^d)$ 上强连续、保正的收缩半群。对 $f\in\mathcal S$ 与 $t>0$，
\[
 \partial_tP_tf=\Delta P_tf=P_t\Delta f.
\]
并且在 $L^p$ 中有右导数
\[
 \lim_{t\downarrow0}\frac{P_tf-f}{t}=\Delta f.
\]
若 $G_p$ 表示该 $C_0$ 半群按定义~\ref{def:4-4-3-2} 得到的生成元，则
\begin{equation}\label{eq:7-2-heat-generator-closure}
 G_p=A_p:=\overline{\Delta|_{\mathcal S}}^{\,L^p\times L^p}.
\end{equation}
换言之，
\[
 D(A_p)=\left\{f\in L^p:\begin{array}{l}
   \text{存在 }\varphi_n\in\mathcal S,
   \ \varphi_n\to f\text{ 于 }L^p,\\[-2pt]
   \Delta\varphi_n\text{ 于 }L^p\text{ 收敛}
 \end{array}\right\},
 \quad A_pf=\lim_n\Delta\varphi_n.
\]
特别地，$\mathcal S$ 是 $A_p$ 的核心。此陈述在 $p=1$ 时不把 $D(A_1)$ 改写成
$W^{2,1}$。
\end{proposition}

\begin{proof}
Gaussian 积分给出 $p_t\geq0$ 与 $\int p_t=1$，故 Young 不等式说明
$\|P_tf\|_p\leq\|f\|_p$，且 $P_t$ 保正。由例~\ref{ex:7-2-2-1}，
\[
 \widehat{p_t}(\xi)=e^{-4\pi^2t|\xi|^2}.
\]
于是卷积公式给出
$\widehat{p_s*p_t}=e^{-4\pi^2(s+t)|\xi|^2}=\widehat{p_{s+t}}$。
定理~\ref{theorem:7-2-2-3} 的函数唯一性应用于二者之差，得到
$p_s*p_t=p_{s+t}$，从而 $P_sP_t=P_{s+t}$。又因
$p_t(x)=t^{-d/2}p_1(x/\sqrt t)$ 是近似恒等族，定理~\ref{theorem:4-4-2-1} 给出
$P_tf\to f$ 于 $L^p$，$1\leq p<\infty$。半群律把零点处强连续性移到每个时刻。

若 $f\in\mathcal S$，则
$e^{-4\pi^2t|\xi|^2}\widehat f(\xi)\in\mathcal S$。在任意闭的正时间区间上可逐个 Schwartz
半范数对乘子求导；定理~\ref{theorem:7-2-3-1} 的反变换给出
\begin{align*}
 \mathcal F(\partial_tP_tf)
 &=-4\pi^2|\xi|^2e^{-4\pi^2t|\xi|^2}\widehat f(\xi),\\
 \mathcal F(\Delta P_tf)
 &=-4\pi^2|\xi|^2e^{-4\pi^2t|\xi|^2}\widehat f(\xi).
\end{align*}
Fourier 的单射性给出热方程。并且
\[
 P_tf-f=\int_0^tP_s\Delta f\dd s
\]
作为 $L^p$ 值积分成立。具体地，对 $0<\eta<t$，正时间上的可微性与微积分基本定理给
\[
 P_tf-P_\eta f=\int_\eta^t\partial_sP_sf\dd s
 =\int_\eta^tP_s\Delta f\dd s.
\]
当 $\eta\downarrow0$ 时，$P_\eta f\to f$；又因 $\Delta f\in\mathcal S\subset L^p$，
$P_s\Delta f\to\Delta f$，故右端在 $L^p$ 中收敛到从 $0$ 开始的 Bochner 积分。于是
\[
 \left\|\frac{P_tf-f}{t}-\Delta f\right\|_p
 \leq\frac1t\int_0^t\|P_s\Delta f-\Delta f\|_p\dd s
 \longrightarrow0.
\]
因此
$\mathcal S\subset D(G_p)$ 且 $G_pf=\Delta f$。

生成元 $G_p$ 是闭算子，故 $\Delta|_{\mathcal S}$ 的图闭包 $A_p$ 满足
$A_p\subset G_p$。反过来取 $f\in D(G_p)$。一般 $C_0$ 半群的生成元交换性质
\ref{proposition:4-4-3-1} 给出
\[
 f_\varepsilon:=P_\varepsilon f\in D(G_p),\qquad
 G_pf_\varepsilon=P_\varepsilon G_pf.
\]
当 $\varepsilon\downarrow0$ 时，$f_\varepsilon\to f$ 且
$G_pf_\varepsilon\to G_pf$ 于 $L^p$。

固定 $\varepsilon>0$。取 $g_n\in\mathcal S$ 使 $g_n\to f$ 于 $L^p$，并令
$\varphi_n=P_\varepsilon g_n$。由 Fourier 乘子表示及 Schwartz 自同构，
$\varphi_n\in\mathcal S$。收缩性给出
\[
 \|\varphi_n-f_\varepsilon\|_p\leq\|g_n-f\|_p\longrightarrow0.
\]
此外
\[
 \partial_t p_t(x)
 =\left(\frac{|x|^2}{4t^2}-\frac d{2t}\right)p_t(x)
 =\Delta p_t(x).
\]
当 $t$ 限于 $[\varepsilon/2,3\varepsilon/2]$ 时，右端及时间差商由一个可积的
“多项式乘 Gaussian”控制；支配收敛给出
$\|(p_{\varepsilon+h}-p_\varepsilon)/h-\Delta p_\varepsilon\|_1\to0$。因此
\[
 G_pf_\varepsilon=(\Delta p_\varepsilon)*f,
 \qquad
 \Delta\varphi_n=(\Delta p_\varepsilon)*g_n.
\]
这两个等式也可由 Gaussian 的显式导数和 Fubini 得到。因此
\[
 \|\Delta\varphi_n-G_pf_\varepsilon\|_p
 \leq\|\Delta p_\varepsilon\|_1\|g_n-f\|_p\longrightarrow0.
\]
取 $\varepsilon_k\downarrow0$，并在第 $k$ 步选 $n(k)$，使相应的
$\varphi_k:=P_{\varepsilon_k}g_{n(k)}$ 满足
\[
 \|\varphi_k-f_{\varepsilon_k}\|_p<\frac1k,
 \qquad
 \|\Delta\varphi_k-G_pf_{\varepsilon_k}\|_p<\frac1k.
\]
于是
\begin{align*}
 \|\varphi_k-f\|_p
 &\leq k^{-1}+\|f_{\varepsilon_k}-f\|_p\longrightarrow0,\\
 \|\Delta\varphi_k-G_pf\|_p
 &\leq k^{-1}+\|G_pf_{\varepsilon_k}-G_pf\|_p\longrightarrow0.
\end{align*}
这给出了图闭包所需的同一条逼近列，并证明
$f\in D(A_p)$、$A_pf=G_pf$，故 $G_p\subset A_p$，并完成式
\eqref{eq:7-2-heat-generator-closure} 的证明。上述稠密、收缩与
$\|\Delta p_\varepsilon\|_1$ 估计在 $p=1$ 仍全部成立，所以端点 $p=1$ 已包含在图闭包论证中。
\end{proof}

生成元结论逐个固定 $p<\infty$ 解释。它既不声称不同 $p$ 的定义域相同，也不声称热半群在
$L^\infty$ 上强连续。具体地，令 $f=\1_{\{x_1>0\}}$，并记 $\Phi$ 为标准一维正态分布函数。
热核的坐标乘积结构给出
\[
 P_tf(x)=\Phi\!\left(\frac{x_1}{\sqrt{2t}}\right).
\]
当 $x_1\downarrow0$ 时右侧趋于 $1/2$，而 $f(x)=1$；每个靠近跳面的正侧薄层都有正测度，故
$\|P_tf-f\|_\infty\geq1/2$ 对每个 $t>0$ 成立。于是该半群在 $L^\infty$ 上不强连续。

\begin{proposition}[Poisson 半群与调和延拓]\label{proposition:7-2-3-4}
对 $t>0$，令
\[
 q_t(x)=c_d\frac{t}{(t^2+|x|^2)^{(d+1)/2}},
 \qquad
 c_d=\frac{\Gamma((d+1)/2)}{\pi^{(d+1)/2}},
 \qquad Q_tf=q_t*f.
\]
则对 $s,t>0$，
\[
 \widehat{q_t}(\xi)=e^{-2\pi t|\xi|},
 \qquad Q_sQ_t=Q_{s+t}.
\]
若 $f\in L^p(\R^d)$，$1\leq p<\infty$，则
$Q_tf\to f$ 于 $L^p$；若 $f\in C_0(\R^d)$，则一致收敛。对这些数据，卷积代表
$u(x,t)=Q_tf(x)$ 在上半空间 $\R^d\times(0,\infty)$ 调和：
\[
 (\Delta_x+\partial_t^2)u=0.
\]
\end{proposition}

\begin{proof}
先证明 Poisson 核与热核的从属关系（subordination）。对 $t>0$、$\lambda\geq0$，
\begin{equation}\label{eq:7-2-subordination-scalar}
 \frac{t}{2\sqrt\pi}\int_0^\infty
 s^{-3/2}e^{-t^2/(4s)}e^{-\lambda s}\dd s=e^{-t\sqrt\lambda}.
\end{equation}
为核对该积分，令 $r=t/(2\sqrt s)$，并记
$a=t\sqrt\lambda/2$。左边化为
\[
 \frac2{\sqrt\pi}J(a),
 \qquad J(a)=\int_0^\infty e^{-r^2-a^2/r^2}\dd r.
\]
当 $a>0$ 时可在积分号下求导；再在所得积分中换元 $u=a/r$，有
\[
 J'(a)=-2a\int_0^\infty r^{-2}e^{-r^2-a^2/r^2}\dd r=-2J(a).
\]
而 $J(0)=\sqrt\pi/2$，故 $J(a)=(\sqrt\pi/2)e^{-2a}$；由单调收敛也覆盖 $a=0$。
这证明式~\eqref{eq:7-2-subordination-scalar}。

将热核代入并使用 Tonelli 定理，得到
\begin{align*}
 &\frac{t}{2\sqrt\pi}\int_0^\infty
  s^{-3/2}e^{-t^2/(4s)}p_s(x)\dd s\\
 &\quad=\frac{t}{2\sqrt\pi}(4\pi)^{-d/2}
  \int_0^\infty s^{-(d+3)/2}
  e^{-(t^2+|x|^2)/(4s)}\dd s\\
 &\quad=\frac{\Gamma((d+1)/2)}{\pi^{(d+1)/2}}
       \frac{t}{(t^2+|x|^2)^{(d+1)/2}}=q_t(x),
\end{align*}
其中最后一步使用 $r=(t^2+|x|^2)/(4s)$。再用 Fubini、热核 Fourier 公式与
式~\eqref{eq:7-2-subordination-scalar}（取 $\lambda=4\pi^2|\xi|^2$），得到
\[
 \widehat{q_t}(\xi)
 =\frac{t}{2\sqrt\pi}\int_0^\infty
 s^{-3/2}e^{-t^2/(4s)}e^{-4\pi^2s|\xi|^2}\dd s
 =e^{-2\pi t|\xi|}.
\]
取 $\xi=0$ 可见 $\int q_t=1$。乘子相乘并使用 $L^1$ Fourier 唯一性，得到
$q_s*q_t=q_{s+t}$，即 Poisson 半群律。

直接核对调和性。置 $\rho=t^2+|x|^2$、$\alpha=(d+1)/2$，暂略常数 $c_d$。对
$q=t\rho^{-\alpha}$，
\begin{align*}
 \Delta_xq
 &=-2\alpha dt\rho^{-\alpha-1}
   +4\alpha(\alpha+1)t|x|^2\rho^{-\alpha-2},\\
 \partial_t^2q
 &=-6\alpha t\rho^{-\alpha-1}
   +4\alpha(\alpha+1)t^3\rho^{-\alpha-2}.
\end{align*}
两式相加后，$2(\alpha+1)=d+3$ 使系数抵消，故
$(\Delta_x+\partial_t^2)q_t=0$。在任意正时间带
$0<a\leq t\leq b$ 上，对实际用到的 $|\gamma|+j\leq3$，直接对
$t(t^2+|x|^2)^{-(d+1)/2}$ 求导并把有界的 $t$ 次幂吸收进常数，可得
\[
 \sup_{a\leq t\leq b}
 \bigl|\partial_t^j\partial_x^\gamma q_t(x)\bigr|
 \leq C_{a,b,\gamma,j}(1+|x|)^{-d-1}.
\]
右端属于每个 $L^r(\R^d)$，$1\leq r\leq\infty$。下面把端点所需的导数交换写清。若
$1<p<\infty$，令 $p'=p/(p-1)$；对任一这样的核导数 $k$，其平移在 $L^{p'}$ 中连续，故差商恒等式
\[
 \frac{k(\,\cdot+h e_j)-k}{h}
 =\int_0^1\partial_jk(\,\cdot+\theta h e_j)\dd\theta
\]
给出差商在 $L^{p'}$ 中趋于 $\partial_jk$，再由 \Holder{} 不等式把极限送入与 $f$ 的卷积。
当 $p=1$ 时，核的下一阶导数连续且在无穷远消失，因而一致连续；同一恒等式给出差商在
$L^\infty$ 中收敛，再与 $f\in L^1$ 配对。对时间差商及二阶导数重复此论证，得到
\[
 \partial_t^2(q_t*f)=(\partial_t^2q_t)*f,
 \qquad \Delta_x(q_t*f)=(\Delta_xq_t)*f.
\]
这些卷积随 $(x,t)$ 连续，因此 $u=Q_tf$ 是经典调和函数。若 $f\in C_0\subset L^\infty$，则改用
核导数在 $L^1$ 中的平移连续性与 $L^1$ 差商收敛，结论相同。

最后，$q_t(x)=t^{-d}q_1(x/t)$，且 $q_1\geq0$、$\int q_1=1$。这是近似恒等族。
定理~\ref{theorem:4-4-2-1} 分别给出 $L^p$（$p<\infty$）强收敛与 $C_0$ 一致收敛。
\end{proof}

这里没有断言 $Q_t$ 在 $L^\infty$ 上强连续。以
$f=\1_{\{x_1>0\}}$ 为例，对其余坐标积分后得到一维 Poisson 核，因而
\[
 Q_tf(x)=\frac12+\frac1\pi\arctan\frac{x_1}{t}.
\]
当 $x_1\downarrow0$ 时右侧趋于 $1/2$，而 $f(x)=1$；正侧任意薄层都有正测度，所以
$\|Q_tf-f\|_\infty\geq1/2$ 对每个 $t>0$ 成立，强连续性确实失败。

\begin{example}[Gaussian 与 Hermite]\label{ex:7-2-3-1}
函数 $g(x)=e^{-\pi|x|^2}$ 属于 $\mathcal S$，且
$\widehat g=g$。命题~\ref{proposition:7-2-2-2} 给出具体的第一阶关系
\[
 \widehat{x_jg}(\xi)=-i\xi_jg(\xi),
 \qquad
 \widehat{\partial_jg}(\xi)=2\pi i\xi_jg(\xi).
\]
反复乘 $x_j$ 与求导，会得到“多项式乘 Gaussian”的函数族；Leibniz 公式说明它们全在
$\mathcal S$。令
\[
 V_N=\{P(x)g(x):\deg P\leq N\}.
\]
对 $|\alpha|\leq N$，乘法--微分公式逐项给出
\[
 \widehat{x^\alpha g}(\xi)
 =\left(\frac{i}{2\pi}\right)^{|\alpha|}\partial^\alpha g(\xi)
 =Q_\alpha(\xi)g(\xi),
 \qquad \deg Q_\alpha\leq|\alpha|.
\]
故 $\mathcal F(V_N)\subset V_N$。反 Fourier 变换等于 Fourier 变换后再作反射，而反射保持
$V_N$，所以实际有 $\mathcal F(V_N)=V_N$。这些 $V_N$ 是按次数递增的有限维子空间。按总次数对
“多项式乘 Gaussian”作正交化可得到 Hermite 函数；又因 $\mathcal F$ 在 $L^2$ 上为酉算子，它保持
$V_N\ominus V_{N-1}$（约定 $V_{-1}=\{0\}$），故每一总次数层内的 Fourier 作用都归结为有限维酉
矩阵。因而 Hermite 函数提供了可显式计算的 Fourier 测试族，并将在谐振子谱问题中再次出现。
\end{example}

\begin{example}[Dirac 分布]\label{ex:7-2-3-2}
定义 $\langle\delta_0,\varphi\rangle=\varphi(0)$。由于
$|\varphi(0)|\leq p_{0,0}(\varphi)$，有 $\delta_0\in\mathcal S'$。按定义
\ref{def:7-2-3-4}，
\[
 \langle\widehat{\delta_0},\varphi\rangle
 =\langle\delta_0,\widehat\varphi\rangle
 =\widehat\varphi(0)=\int\varphi(x)\dd x
 =\langle1,\varphi\rangle.
\]
所以 $\widehat{\delta_0}=1$。反过来，Schwartz 反演给出
\[
 \langle\widehat1,\varphi\rangle
 =\int\widehat\varphi(x)\dd x=\varphi(0)
 =\langle\delta_0,\varphi\rangle,
\]
故 $\widehat1=\delta_0$。这些等式只表示两个缓增分布对每个 Schwartz 测试函数的作用相同，并未把
$\delta_0$ 当作普通点值函数。
\end{example}

\begin{example}[热核]\label{ex:7-2-3-3}
把例~\ref{ex:7-2-2-1} 中的参数取为 $a=(4\pi t)^{-1}$，有
\begin{align*}
 \widehat{p_t}(\xi)
 &=(4\pi t)^{-d/2}
   (4\pi t)^{d/2}e^{-4\pi^2t|\xi|^2}\\
 &=e^{-4\pi^2t|\xi|^2}.
\end{align*}
因此 $p_s*p_t=p_{s+t}$ 在频率侧只是
$e^{-4\pi^2s|\xi|^2}e^{-4\pi^2t|\xi|^2}
=e^{-4\pi^2(s+t)|\xi|^2}$。同一个乘子求时间导数给出
$-4\pi^2|\xi|^2$，正好是 $\Delta$ 的 Fourier 符号。
\end{example}

\begin{example}[Brownian 与热半群的同一计算]\label{ex:7-2-3-4}
若第六章的标准 Brownian motion 满足 $B_t\sim N(0,tI)$，则其密度是
\[
 (2\pi t)^{-d/2}e^{-|y|^2/(2t)}=p_{t/2}(y).
\]
所以对有界 Borel 函数 $f$，以及对任意 $f\in L^p(\R^d)$、$1\le p\le\infty$
（此时取卷积所定义的代表元），
\[
 \E f(x+B_t)=\int f(x+y)p_{t/2}(y)\dd y=P_{t/2}f(x),
\]
其中最后一步使用 $p_{t/2}(y)=p_{t/2}(-y)$。若 $f\in\mathcal S$，写
$B_t=\sqrt t\,Z$，其中 $Z\sim N(0,I_d)$。二阶 Taylor 公式给
\[
 f(x+y)-f(x)
 =\nabla f(x)\cdot y+\frac12y^{\mathsf T}D^2f(x)y+r(x,y),
 \qquad
 |r(x,y)|\leq\frac12|y|^2\omega_{D^2f}(|y|),
\]
其中 $\omega_{D^2f}(s)\to0$ 当 $s\downarrow0$。利用
$\E Z_j=0$、$\E Z_jZ_k=\delta_{jk}$，得到
\[
 \sup_x\left|
   \frac{\E f(x+B_t)-f(x)}t-\frac12\Delta f(x)
 \right|
 \leq\frac12\E\!\left[|Z|^2
        \omega_{D^2f}(\sqrt t|Z|)\right]\longrightarrow0.
\]
最后一步使用支配收敛，因为
$\omega_{D^2f}\leq2\|D^2f\|_\infty$ 且 $\E|Z|^2<\infty$。特别地，
\[
 \lim_{t\downarrow0}\frac{\E f(x+B_t)-f(x)}t
 =\frac12\Delta f(x).
\]
因此概率时间 $t$ 对应本节 PDE 热时间 $t/2$：Brownian 生成元是
$\tfrac12\Delta$，而 $P_t$ 的生成元是 $\Delta$。方差参数中的这个 $1/2$ 不能由记号省略。
\end{example}

图~\ref{fig:7-2-schwartz-plancherel-semigroups} 汇总本小节的三层结构。内层的 Schwartz 方块允许
逐点积分、微分与反演；垂直的稠密完备化把它唯一扩张为 $L^2$ 上的酉算子；右侧两个指数乘子则把连续
时间演化化成乘法。生成元属于无界算子层，必须连同图闭包定义域理解。

\begin{figure}[htbp]
\centering
\begin{tikzpicture}[node distance=1.25cm and 2.0cm]
  \node[book strong node] (sx) {$\mathcal S_x$};
  \node[book strong node,right=of sx] (sxi) {$\mathcal S_\xi$};
  \draw[book accent arrow,bend left=14] (sx) to node[above,book label]{$\mathcal F$} (sxi);
  \draw[book accent arrow,bend left=14] (sxi) to node[below,book label]{$\mathcal F^{-1}$} (sx);

  \node[book node,below=1.55cm of sx] (l2x) {$L^2_x$};
  \node[book node,below=1.55cm of sxi] (l2xi) {$L^2_\xi$};
  \draw[book dashed arrow] (sx)--node[left,book label]{稠密完备化} (l2x);
  \draw[book dashed arrow] (sxi)--node[right,book label]{稠密完备化} (l2xi);
  \draw[book arrow,bend left=12] (l2x) to node[above,book label]{酉} (l2xi);
  \draw[book arrow,bend left=12] (l2xi) to (l2x);

  \node[book warm node,right=2.0cm of sxi] (heat)
    {热乘子\\$e^{-4\pi^2t|\xi|^2}$};
  \node[book warm node,below=1.2cm of heat] (pois)
    {Poisson 乘子\\$e^{-2\pi t|\xi|}$};
  \draw[book arrow] (sxi)--(heat);
  \draw[book arrow] (sxi)--(pois);
  \node[book warning node,below=1.0cm of pois] (gen)
    {生成元\\图闭包与定义域};
  \draw[book dashed arrow] (heat)--(gen);
\end{tikzpicture}
\caption{Schwartz 自同构、$L^2$ 完备化与热/Poisson 乘子。}
\label{fig:7-2-schwartz-plancherel-semigroups}
\end{figure}

\begin{exercise}
\begin{enumerate}
  \item 用 Leibniz 公式逐项证明 $\mathcal S(\R^d)$ 对微分和乘多项式闭合；再对
        $f,g\in\mathcal S$ 把积分分成 $|y|\le |x|/2$ 与其补集，证明 $f*g\in\mathcal S$。
  \item 从恒等式
        $\xi^\alpha\partial_\xi^\beta\widehat f
        =(2\pi i)^{-|\alpha|}(-2\pi i)^{|\beta|}
        \widehat{\partial^\alpha(x^\beta f)}$
        出发，对每个输出半范数写出有限个输入半范数的显式控制，从而证明
        $\mathcal F:\mathcal S\to\mathcal S$ 连续。
  \item 只用 $\mathcal S$ 在 $L^2$ 中稠密、Parseval 等式与 $L^2$ 完备性，构造 Fourier 的
        $L^2$ 延拓；逐步证明定义与逼近列无关，并用反 Fourier 延拓证明其像为整个 $L^2$。
  \item 从乘子 $e^{-4\pi^2t|\xi|^2}$ 推出热算子的半群律。对 $f\in\mathcal S$ 证明乘子可在
        $t$ 上求导并得到 $\partial_tP_tf=\Delta P_tf$；再指出近似恒等性质为何只给
        $1\le p<\infty$ 上的强连续性，而不自动给 $L^\infty$ 情形。
\end{enumerate}
\end{exercise}

圆周与欧氏空间的 Fourier 理论都从所有指数测试函数上的积分值唯一恢复对象：圆周用 \Fejer{} 正核和离散
$\ell^2$ 完备化，欧氏空间用 Gaussian 正则化和连续 $L^2$ 完备化。下一节将不再预先给出测度，
而从一列候选矩出发，研究正定性与紧性何时足以保证代表测度存在。
